\documentclass[11pt]{article}
\usepackage[utf8]{inputenc}
\usepackage{amsmath,amssymb}
\usepackage{graphicx}
\usepackage{booktabs}
\usepackage{authblk}
\usepackage[margin=1in]{geometry}
\usepackage[hidelinks]{hyperref}
\usepackage{xurl}
\usepackage{setspace}

\title{Magnetism as Circulating Charge:\\
Why the Nuclear Electron Carries No Magnetic Moment}
\author[1]{Claes Johnson}
\affil[1]{KTH Royal Institute of Technology, SE-100 44 Stockholm, Sweden \\
\texttt{claesjohnson@gmail.com}}
\date{\today \\[0.4em]
\small\textit{Argument and claims by the author; drafting and computation prepared with AI assistance (Claude, Anthropic), reviewed and verified by the author.}}

\begin{document}
\maketitle

\begin{abstract}
The classic objection that retired the proton--electron model of the nucleus in 1932 is the
\emph{magnetic-moment} problem: an electron confined inside a nucleus should carry a magnetic moment of
order the Bohr magneton $\mu_B$, whereas measured nuclear moments --- the neutron's included --- are of
order the nuclear magneton $\mu_N \approx \mu_B/1836$, a thousand-fold too small. We show that in the
RealQM/RealNucleus picture, where constituents are charge densities rather than point particles, this
objection dissolves. In a charge-density theory the magnetic moment is a property of the \emph{current},
$\boldsymbol\mu = \tfrac12\!\int \mathbf r\times\mathbf J\,d^3r$ with
$\mathbf J = (\hbar/m)\,\mathrm{Im}(\psi^*\nabla\psi)$, and a \emph{real} density carries no current and
hence no moment. RealNucleus computes the deuteron's confined electron as a \emph{flat, real,
node-less} charge density; it therefore carries \emph{exactly zero} magnetic moment, and the nuclear
moment is set by the protons alone --- at $\mu_N$-scale, not $\mu_B$-scale. A direct numerical test
confirms the enabling facts: a single global phase winding $e^{im\phi}$ gives a moment $\mu_z=-m\,\mu_B$
that is \emph{quantized and independent of the density's spatial extent} (so caging cannot continuously
suppress a winding), while only the winding-free ($m=0$) real density gives zero. The same flatness that
was shown to make the electron's \emph{mass} irrelevant to nuclear binding thus also makes its
\emph{magnetic moment} vanish --- one property answering two of the classical objections at once. We are
explicit about what remains open: the free electron's own $\mu_B$ and the neutron's exact
$-1.913\,\mu_N$ require a structured, self-consistent current and hence a complex/real-time solver, and
the anomalous $g/2 = 1.00116\ldots$ is a separate debt.

\medskip
\noindent\textbf{Keywords:} nuclear magnetic moment; nuclear electrons; RealQM; RealNucleus; circulating
charge; charge-density current; flat electron; neutron.
\end{abstract}

\setstretch{1.12}

\section{The objection}
\label{sec:objection}

When the proton--electron model of the nucleus was abandoned in 1932, two arguments did the work. One
was the confinement energy (an electron localised to nuclear size should carry hundreds of MeV of kinetic
energy); the other was the \emph{magnetic moment}. The magnetic moment argument is the sharper of the
two. An electron carries an intrinsic magnetic moment of about one Bohr magneton,
\begin{equation}
\mu_B = \frac{e\hbar}{2m_e} \;=\; 1836\,\mu_N, \qquad \mu_N = \frac{e\hbar}{2m_p},
\end{equation}
so any nucleus containing electrons should carry moments of order $\mu_B$. Measured nuclear moments are
instead of order $\mu_N$: the neutron's is $\mu_n = -1.913\,\mu_N \approx -0.001\,\mu_B$. An electron
inside would overshoot by a factor of roughly a thousand. In the standard picture the difficulty never
arises, because the neutron is not $p+e$ but an elementary $udd$ baryon whose moment comes from its
charged quarks.

RealNucleus revives the proton--electron nucleus \cite{RealNucleus}: the neutron is a bound $p+e$ pair,
the deuteron is $2p+1e$, and light nuclei are dual-Coulomb-confined arrangements of proton and electron
charge densities. If that picture is to survive, it must answer the magnetic-moment objection. This note
shows that it does --- and that the answer is the \emph{same} flat electron that answers the
confinement-energy objection.

\section{Magnetic moment as a property of the current}
\label{sec:current}

The decisive difference between a point-particle theory and a charge-density theory is where the magnetic
moment comes from. For a point electron the dominant moment is \emph{intrinsic} spin --- a fixed
$\sim\mu_B$ that the particle carries into any state, and which cannot be switched off by confinement.
In a charge-density theory there is no such term \emph{a priori}; the magnetic moment is generated
entirely by the \emph{current} in the density,
\begin{equation}
\boldsymbol\mu \;=\; \frac{1}{2}\!\int \mathbf r\times\mathbf J\,d^3r,
\qquad
\mathbf J \;=\; \frac{\hbar}{m}\,\mathrm{Im}\!\left(\psi^*\nabla\psi\right).
\label{eq:mu}
\end{equation}
The current is the electromagnetic reading of charge in motion, and it has one immediate and decisive
property: for a \emph{real} wavefunction $\psi$, the combination $\psi^*\nabla\psi$ is real, so
$\mathbf J = 0$ identically. \textbf{A static, real charge density carries no current and therefore no
magnetic moment.} Magnetism, in this framework, requires the complex, time-dependent structure of the
wavefunction --- specifically a phase that \emph{winds in space}, which is the charge-density realisation
of ``charge moving in circles.''

A phase winding $\psi = R(\mathbf r)\,e^{im\phi}$ about an axis produces a genuine azimuthal current and,
through \eqref{eq:mu}, an orbital magnetic moment. Its size is the subject of the next section.

\section{The winding is quantised: a numerical test}
\label{sec:test}

Impose a phase winding $e^{im\phi}$ on a normalised electron density and evaluate \eqref{eq:mu}. Writing
$s$ for the cylindrical radius, the current is $J_\phi = (\hbar/m)\,|R|^2\,(m/s)$, so
$(\mathbf r\times\mathbf J)_z = s\,J_\phi = (\hbar/m)\,m\,|R|^2$, and
\begin{equation}
\mu_z \;=\; -\frac{e}{2}\!\int (\hbar/m)\,m\,|R|^2\,d^3r \;=\; -\,m\,\mu_B ,
\label{eq:quantised}
\end{equation}
since $|R|^2$ is normalised. \textbf{The moment is $-m\,\mu_B$ exactly, independent of the shape or
extent of $|R|^2$}: the factors of $s$ and $1/s$ cancel. A direct grid evaluation confirms this and
isolates the escape route (Table~\ref{tab:test}).

\begin{table}[h]
\centering
\small
\renewcommand{\arraystretch}{1.15}
\begin{tabular}{lll}
\toprule
configuration & current pattern & $\mu_z/\mu_B$ \\
\midrule
free electron (broad, $\sigma=5$)     & single winding $m=1$ & $-1.0000$ \\
caged electron (narrow, $\sigma=0.5$) & single winding $m=1$ & $-1.0000$ \\
caged electron (very narrow, $\sigma=0.2$) & single winding $m=1$ & $-1.0000$ \\
any real density                      & none ($m=0$)         & $\phantom{-}0.0000$ \\
\midrule
caged, core$+$/halo$-$ current        & structured (counter-rotating) & $+0.212$ \\
\bottomrule
\end{tabular}
\caption{Magnetic moment from imposed currents on normalised densities. A single global winding gives
$-m\,\mu_B$ regardless of the density's spatial extent (rows 1--3): geometry does \emph{not} suppress a
winding. Only the winding-free real density gives zero. A \emph{structured} current with opposed
circulation in core and halo gives a non-quantised, tunable moment --- the only route to a small nonzero
value.}
\label{tab:test}
\end{table}

Two conclusions follow, one negative and one positive. \emph{Negative:} caging cannot
\emph{continuously} suppress a magnetic moment --- a single winding yields the full $\mu_B$ whether the
electron is spread over five length units or squeezed into one fifth of one. The naive picture of
``confinement shrinks the moment'' is simply wrong. \emph{Positive:} the only winding that gives zero is
$m=0$, the winding-free \emph{real} density.

\section{Energetics: why the caged electron cannot afford a winding}
\label{sec:energetics}

Section~\ref{sec:test} shows that geometry does not \emph{continuously} shrink a winding --- a winding of
$m=1$ gives $\mu_B$ at any size. But there remains the prior question of whether a caged electron can
\emph{sustain} a winding at all, and this is settled by energetics. A winding carries a kinetic cost,
\begin{equation}
\Delta T(m) \;=\; \frac{\hbar^2}{2m_e}\,m^2\!\int \frac{|R|^2}{s^2}\,d^3r
\;=\; \mathrm{Ry}\;m^2\,(a_0/\sigma)^2 ,
\label{eq:cost}
\end{equation}
where $\sigma$ is the density's cylindrical scale and the last form uses the $m=1$ node-on-axis profile
$|R|^2\propto s^2 e^{-s^2/\sigma^2}$ (for which $\langle1/s^2\rangle = 1/\sigma^2$ exactly), with
$\mathrm{Ry}=\hbar^2/2m_e a_0^2 = 13.6$~eV. The cost scales as $1/\sigma^2$, so it is minute for a
spread-out free electron and enormous for a compact caged one (Table~\ref{tab:cost}).

\begin{table}[h]
\centering
\small
\renewcommand{\arraystretch}{1.15}
\begin{tabular}{lll}
\toprule
electron & scale $\sigma$ & winding cost $\Delta T(m{=}1)$ \\
\midrule
free (Bohr scale)          & ${\sim}a_0$ & $13.6$~eV \quad($\approx 1$~Ry) \\
caged (deuteron)           & ${\sim}2$~fm & $9.5$~GeV \\
caged (deuteron)           & ${\sim}1$~fm & $38$~GeV \\
\bottomrule
\end{tabular}
\caption{Kinetic cost of an $m{=}1$ winding, Eq.~\eqref{eq:cost}. For the free electron it is about one
Rydberg --- affordable, so the free electron can sustain a circulation. For the femtometre-caged electron
it is of order GeV, thousands of times the ${\sim}2.2$~MeV deuteron binding --- so a winding is forbidden
and the electron collapses to the winding-free $m{=}0$ state. A winding is a gradient; it reintroduces
exactly the kinetic cost the flat electron avoids.}
\label{tab:cost}
\end{table}

The conclusion is robust, though the exact margin depends on a mass convention. Using the physical
electron mass, the caged winding ($\sim\!9.5$~GeV at $2$~fm) exceeds the deuteron binding by a factor
$\sim\!10^{3}$. Using RealNucleus's equal-mass device ($m_e\!\to\!m_p$ in the kinetic term) reduces it by
$1836$ to $\sim\!5$~MeV --- still above the $2.2$~MeV binding, so still forbidden, but only by a factor
$\sim\!2$. Either way the caged electron cannot afford the winding and settles at $m=0$; only the size of
the margin is convention-dependent. Crucially this is the \emph{same} statement as mass-irrelevance: the
flat $m=0$ electron avoids femtometre-scale gradients, and a winding is precisely such a gradient, so the
property that removes the mass also forbids the current.

\section{The flat electron carries no moment}
\label{sec:flat}

This is where the RealNucleus computation supplies the answer. The deuteron's confined electron is not an
arbitrary state: minimisation of the Coulomb energy under the free-boundary condition returns it as a
\emph{flat, node-less, real} charge density, caged by the two protons and satisfying charge continuity at
a definite radius \cite{RealNucleus}. A flat real density is precisely the $m=0$ case: $\psi$ real
$\Rightarrow \mathbf J = 0 \Rightarrow \boldsymbol\mu = 0$. \textbf{The nuclear electron carries no
magnetic moment.}

The consequence is exactly what is observed. With the confined electron contributing nothing, the
magnetic moment of a nucleus is set by its \emph{protons}, and therefore comes out at
\emph{nuclear-magneton scale}, not Bohr-magneton scale. The factor-of-a-thousand overshoot that retired
the proton--electron model does not occur, because the object that would have supplied it --- an electron
carrying its intrinsic $\mu_B$ into the nucleus --- does not exist in a charge-density theory. There is
no intrinsic term to carry; there is only the current, and the flat electron's current is zero.

\paragraph{One flatness, two objections.} The flat confined electron was not introduced here to solve the
magnetic-moment problem. It was already required, in the companion work \cite{RealNucleus}, to make the
electron's \emph{mass} irrelevant to nuclear binding: a flat density has vanishing kinetic-energy
gradient, so its mass drops out, and the deuteron's femtometre scale is set by the heavy proton rather
than by confining a light electron at hundreds of MeV. That same flatness --- real, node-less,
currentless --- now answers the magnetic-moment objection as well. The two classical arguments that
retired nuclear electrons in 1932, the confinement energy and the magnetic moment, are answered by a
single property of the confined-electron solution. This is not two patches but one fact read twice.

\section{What remains open}
\label{sec:open}

We state plainly what this does and does not establish.

\emph{Established:} the confined electron, computed as a flat real density, carries no magnetic moment, so
RealNucleus predicts nuclear moments at $\mu_N$-scale rather than $\mu_B$-scale --- the qualitative heart
of the objection is removed, and by the same flatness that handles the mass.

\emph{Open, and requiring a complex/real-time solver:}
\begin{itemize}
\item \textbf{The free electron's own moment --- and why it is not an ordinary winding.} A free electron
carries $\sim\mu_B$, and it is tempting to attribute this to an orbital winding of the kind used above.
But that will not do, and the distinction is important. The ground state of a free or atomically-bound
electron is \emph{real} (the $1s$ state, $m=0$): it has \emph{zero orbital moment}, and in the standard
account its $\mu_B$ is \emph{intrinsic spin}, not orbital circulation. So the free electron's $\mu_B$
cannot come from an atomic-scale winding --- if it circulates at all, it must be an \emph{internal},
Compton-scale circulation ($s\sim\hbar/m_ec$), the Zitterbewegung structure, which is a far deeper and
more speculative object than the atomic-scale current considered here. This means the clean, defensible
result of this note is the \emph{caged} side --- the confined electron has no moment --- while the
\emph{free} electron's $\mu_B$ is a genuinely open problem at the Compton scale, not settled by any
atomic-scale winding calculation. We are careful not to conflate the two: suppressing the caged moment is
solid; generating the free $\mu_B$ from circulation is a program.
\item \textbf{The neutron's exact $-1.913\,\mu_N$.} This is small but nonzero, so it is not ``electron
contributes exactly nothing'' but a proton-dominated value with a small residual. Table~\ref{tab:test}
shows a nonzero, non-quantised moment requires a \emph{structured} current (opposed circulation in core
and halo), which only a self-consistent complex solution can generate --- not an imposed single winding.
\item \textbf{The anomalous moment.} The precisely measured $g/2 = 1.00116\ldots$ is QED's crown jewel; a
circulating-charge account would owe its own derivation of this correction. This is the hardest debt and
is untouched here.
\end{itemize}

None of these is hidden, and none is answered by hand-waving. What is claimed is bounded and, we believe,
correct: the \emph{scale} objection --- the thousand-fold overshoot that is the sharp form of the
magnetic-moment argument against nuclear electrons --- is removed, because a flat real confined electron
has no current and hence no moment.

\section{Conclusion}
\label{sec:conclusion}

In a point-particle theory an electron's magnetic moment is intrinsic and cannot be switched off, so a
nuclear electron wrecks the moment budget by a factor of a thousand. In a charge-density theory the
moment is the current's, a real density has no current, and the confined electron --- computed flat and
real --- has no moment. Nuclear moments then sit at $\mu_N$-scale, as observed, and the same flatness
that makes the electron's mass irrelevant to binding makes its magnetic moment vanish. The magnetic
moment, long the sharpest weapon against nuclear electrons, turns out to depend on treating the electron
as a point that carries spin; read as a charge density that carries current, a caged flat electron
carries nothing. What remains --- the free electron's $\mu_B$, the neutron's exact value, the anomalous
$g$-factor --- is well-posed and is the province of a complex, current-bearing extension of the solver.

\begin{thebibliography}{9}
\bibitem{RealNucleus} C.~Johnson, \emph{RealNucleus: Nuclear Binding as Dual Coulomb Confinement, without
a Strong or Weak Force}, 2026 (submitted to Physics Essays);
\url{https://claes542.github.io/RealMolecule/gallery.html}.
\bibitem{RealQM} C.~Johnson, \emph{Real Quantum Mechanics as Multiphase 3D Continuum Mechanics},
manuscript, 2026.
\end{thebibliography}

\end{document}
